\documentclass[12pt]{article}
\usepackage[utf8]{inputenc}
\usepackage[T1]{fontenc}
\usepackage{hyperref}
\title{Urarina Language Resources for K-12 Curriculum Development}
\author{Compiled by Claude Sonnet 5 (claude-sonnet-5) for Andr\'as Kornai}
\date{}
\begin{document}
\maketitle

\section{Introduction}
Urarina (ISO 639-3: ura) is a language isolate of Loreto, Peru. It is spoken mainly in the Chambira, Corrientes and Urituyacu basins. One source gives about 3,000 speakers and another (a UNESCO atlas entry) 1,301. Contact with oil companies and traders has led to a shift to Spanish in some communities.

\section{Primary Educational Resources}
\begin{itemize}
\item A Peruvian Ministry of Education school-library file lists a ``Diccionario visual-URARINA'' among materials for intercultural bilingual schools. The extracted text was fragmentary, so its contents are unverified: \url{https://bibliotecaescolar.minedu.gob.pe/bitstreams/e0432f95-fb42-4484-a6de-26d6fc5e1a5d/download}
\end{itemize}

\section{Linguistic and Reference Materials}
\begin{itemize}
\item Olawsky, article on Upper Chambira Urarina, \emph{Journal of the International Phonetic Association} (2019). It mentions an official alphabet and the creation of educational materials: \url{https://resolve.cambridge.org/core/product/3556B5DDB589F632E3F6F4DB5BA1F1B9/core-reader}
\item Culture in Crisis project page, a multilingual Urarina-Spanish-English dictionary and grammar project: \url{https://cultureincrisis.org/projects/language-and-culture-of-the-urarina-people-of-peru-preparation-of-grammar-and-dictionary-for-an-endangered-language}
\item James Cook University research record: \url{https://researchonline.jcu.edu.au/80355/}
\item Wikipedia: \url{https://en.wikipedia.org/wiki/Urarina}
\end{itemize}

\section{Religious and Community Resources}
No scripture or community organisation resources were found in the search results.

\section{Digital and Audio Resources}
\begin{itemize}
\item UNESCO language atlas entry: \url{https://ar.wal.unesco.org/en/node/72622}
\end{itemize}

\section{Curriculum Development Context}
A reported teacher-training programme for Urarina teachers in most Chambira villages, focused on Urarina as the language of education, suggests some bilingual-education infrastructure. This comes from Olawsky's 2023 chapter, which was not directly seen. The Ministry's visual dictionary and the dictionary project should be checked for availability and scope.
\end{document}